\section*{H.1 复现实务说明一：环境配置}
为保证数值实验可重复，建议固定编译器版本、数学库版本和随机种子，并在每次运行时记录主机信息、执行时间戳、配置快照与脚本哈希值。对于多机并行任务，需要保存任务切分方式与节点间通信日志，避免因分配差异导致结果偏移。
\clearpage

\section*{H.2 复现实务说明二：数据管理}
观测数据应按“原始层、清洗层、建模层”分层保存。原始层禁止覆盖；清洗层应记录插值、去噪和异常值处理规则；建模层应保存用于本次实验的最终快照。建议统一以CSV或Parquet形式保存，并附字段字典和单位说明。
\clearpage

\section*{H.3 复现实务说明三：参数版本}
参数管理建议采用版本化策略。每次修改参数边界、初值池或阻尼规则，都应生成独立版本号并记录变更原因。若实验结果出现异常，可快速回滚到上一稳定版本，定位问题来源。
\clearpage

\section*{H.4 复现实务说明四：日志规范}
日志建议分为运行日志、诊断日志与结果日志。运行日志记录任务生命周期；诊断日志记录告警与回退动作；结果日志记录最终指标。三类日志应使用统一时间基准，并支持按任务ID检索。
\clearpage

\section*{H.5 复现实务说明五：性能评估}
性能评估不应只看总耗时，还需关注每步耗时分解，如导数计算、线性方程求解、I/O写入等。通过分解耗时可定位瓶颈并实施定向优化，例如缓存雅可比结构、减少重复内存分配、并行化残差评估。
\clearpage

\section*{H.6 复现实务说明六：误差解释}
误差解释应区分数值误差、建模误差与数据误差。数值误差来自积分截断与迭代容差；建模误差来自方程简化；数据误差来自测量噪声与采样偏差。分层解释有助于提出有针对性的改进方案。
\clearpage

\section*{H.7 复现实务说明七：报告模板}
建议报告统一包含实验目的、输入条件、核心参数、算法配置、关键结果、失败样本、风险评估与后续计划。统一模板可以显著降低团队协作成本，也便于跨实验横向对比。
\clearpage

\section*{H.8 复现实务说明八：上线前核查}
上线前应完成三项核查：结果一致性核查、边界条件核查、异常恢复核查。通过“正常场景 + 极端场景 + 故障场景”三类用例验证后，再进入生产运行，可有效降低线上不稳定风险。
